
Temporal window duration is an important design consideration in EEG-based driving-state monitoring because it can influence both classification performance and system responsiveness~\citep{rashid2021random}. Previous drowsiness and driving-fatigue studies have compared multiple analysis windows and shown that performance can vary with temporal context~\citep{lee2023lstm,rashid2021random}.
In contrast, EEG eye-blink detectors have generally been evaluated using a single predefined epoch duration rather than systematically across multiple durations ~\citep{chang2016detection,wang2025sliding,shahbakhti2021vme}.
A systematic assessment of how epoch duration affects blink-detection performance, particularly in driving EEG, is therefore still lacking.
In the present study, blink-detection performance at the frontopolar electrodes remained largely stable across 10--40~s epochs.
This differs from drowsiness classification, where performance has been shown to vary more noticeably with analysis-window duration
~\citep{lee2023lstm}. One possible explanation is that blink artifacts typically last only about 100--500~ms~\citep{chang2016detection}, substantially shorter than the epochs evaluated here. Once an epoch is sufficiently long to contain the complete blink waveform, further increases in duration may therefore have limited influence on blink detection itself, although they can alter the amount of background EEG contributing to threshold estimation. From a practical perspective, the observed stability suggests that the proposed pipeline can accommodate a range of epoch durations used in driving-EEG protocols without substantial loss of detection performance.
